\section{Related Work}
\label{sec:related}

\todo{Write four subsections and a comparison table. Verify every citation in \texttt{refs.bib}
before submission.}

\subsection{Testing database engines for logic bugs}
Pivoted query synthesis \cite{rigger2020pqs}, non-optimizing reference engines
\cite{rigger2020norec}, query partitioning \cite{rigger2020tlp}, and the SQLancer tool family
\cite{rigger2020sqlancer}. These provide the oracle vocabulary we adopt, but target SQL DBMSs.

\subsection{Differential and metamorphic testing of data systems}
NEZHA \cite{rui2017nezha} for domain-independent differential testing; SparkFuzz
\cite{sparkfuzz2020} for query-engine correctness regressions; BigFuzz \cite{bigfuzz2020} for
data-analytics framework abstraction; Csmith \cite{yang2011csmith} as the classical
differential-fuzzing reference. \todo{Add metamorphic DB testing (e.g., QTRAN) and recent
equivalent-expression / constraint-based oracles with verified metadata.}

\subsection{DataFrame and workload testing}
FuzzyData and DataFrame workload generation/replay. \todo{Verify metadata; position as workload
generation + replay vs our cross-model semantic oracle.}

\subsection{Search, scheduling, and diversity}
Multi-scope adaptive selection, mutation scheduling (MOPT \cite{mopt2019}), and quality-diversity
archives (MAP-Elites \cite{mouret2015mapelites}). \todo{Position as borrowed mechanisms with
new scope unification, not as primary novelty.}

\begin{table}[t]
\caption{Positioning. \todo{Complete from verified related work:
system, target models, oracle families, evidence pipeline, cross-model reuse.}}
\label{tab:related}
\small
\begin{tabular}{@{}lcccc@{}}
\toprule
System & Models & Oracles & Evidence pipeline & Cross-model reuse \\
\midrule
PQS/NoREC/TLP & SQL & 1 & partial & no \\
FuzzyData & DataFrame & \todo{} & \todo{} & no \\
SparkFuzz & query engine & differential & no & no \\
\sys{} & DataFrame+Arrow+SQL+QE & multi-endpoint & yes & yes \\
\bottomrule
\end{tabular}
\end{table}
